%% Generated from paper/draft.md by `make overleaf` (pandoc + paper/copernicus/).
%% Edit draft.md, not this file: it is overwritten on every build.
%% Copernicus Publications LaTeX package 7.16, The Cryosphere (tc), manuscript (discussion) layout.
\documentclass[tc, manuscript]{copernicus}

%% Helpers for pandoc output
\providecommand{\tightlist}{\setlength{\itemsep}{0pt}\setlength{\parskip}{0pt}}

\begin{document}

\title{Climate-aware sensor placement for Greenland Ice Sheet surface
melt (working title)}

%% TODO: authors and affiliations. Add `authors` and `affiliations` to the YAML header of draft.md:
%%   authors:      [{given: Ana, surname: Smith, affil: 1, email: ana@example.org}]
%%   affiliations: [{id: 1, address: "Institute, City, Country"}]
\Author[1][TODO@example.org]{TODO}{TODO}
\affil[1]{TODO}

\runningtitle{TODO}

\runningauthor{TODO}

\received{}
\pubdiscuss{}
\revised{}
\accepted{}
\published{}

\firstpage{1}

\maketitle

\begin{abstract}
TODO. Write last. Suggested structure: motivation (sparse, incrementally
built in-situ network), approach (probabilistic ML emulator of HIRHAM5
melt used to rank sensor locations), main results (emulator skill under
present and future climates; where placements change), implication.
\end{abstract}

\begin{quote}
Draft status: scaffold. Statements about the project design are taken
from the planning document and the project discussions. \texttt{TODO}
marks missing content. Numbers come from
\texttt{results/evaluation\_metrics.csv}.
\end{quote}

\introduction

The Greenland Ice Sheet is losing mass, and surface melt is a major
driver of that loss. TODO: cite. Existing in-situ networks such as
PROMICE and GC-Net were built incrementally and are not optimised for
detecting tipping-point early-warning signals. TODO: cite and confirm
wording.

Probabilistic machine-learning models can suggest informative sensor
placements by finding sites that most reduce prediction uncertainty.
Andersson et al. \citeyearpar{andersson2023} showed this for Antarctic
surface air temperature with a convolutional Gaussian neural process,
assessing placements with a simulated experiment that keeps the search
period separate from the analysis period.

A placement that is optimal for today's climate need not remain so as
the ice sheet warms. Machine-learning emulators trained on one climate
regime degrade under another \citep{schlager_preprint}, so a placement
study across climate windows has to account for emulator
transferability.

We address two questions:

\begin{itemize}
\tightlist
\item
  \textbf{Q1.} Where are the most informative sensor locations under
  present-day climate conditions?
\item
  \textbf{Q2.} Where are they under future climate conditions, and which
  locations are no-regret (informative in both) versus future-only?
\end{itemize}

TODO: contributions paragraph.

\section{Previous work and gaps}\label{previous-work-and-gaps}

\begin{quote}
TODO: the summaries below come from abstracts and publisher records;
read the full texts of \citet{hakim2020}, \citet{tardif2022} and
\citet{eksen2025} before keeping specific numbers. ``To our knowledge''
claims rest on a web search of November 2026.
\end{quote}

Methods for designing polar observing networks fall into two families.
Traditional methods use a physical model and its error statistics.
Ensemble sensitivity ranks candidate sites by how much an observation
there would reduce the total variance of a target field. It was used to
place automatic weather stations (AWS) for Antarctic surface temperature
\citep{hakim2020}. Its predictions agreed with data-assimilation
experiments, and optimal networks beat random ones by up to about 50\%
\citep{tardif2022}. It has also been used to place Arctic Ocean buoys
for 2-m temperature \citep{kim2024}. Observing system simulation
experiments (OSSEs) and quantitative network design evaluate
hypothetical networks inside a data-assimilation system. They have been
applied mainly to Arctic sea ice \citep{kaminski2015, kaminski2018} and
advocated for the Arctic surface network \citep{masutani2013}.
Representativeness-based designs cluster the domain on environmental
covariates and add sites to the least-represented clusters
\citep{hoffman2013, pallandt2022}. Data-driven methods instead rank
sites with a statistical or machine-learning model. Examples are greedy
mutual-information placement with Gaussian processes \citep{krause2008},
neural-process placement for Antarctic temperature
\citep{andersson2023}, acquisition on epistemic uncertainty with
ConvCNPs \citep{eksen2025}, and greedy placement of borehole
thermometers at Greenland and Antarctic ice-core sites
\citep{shaju2025}.

Against this background, we identify four gaps.

\begin{enumerate}
\def\labelenumi{\arabic{enumi}.}
\tightlist
\item
  \textbf{No optimised network for the Greenland Ice Sheet.} To our
  knowledge, AWS placement on the ice sheet has not been optimised.
  PROMICE sites were chosen to complement existing stations within reach
  of airports and heliports \citep{fausto2021}, in slow-moving regions
  considered relatively safe from crevasses \citep{fausto2026}. The
  design of a Greenland observing system remains an open question
  \citep{grioos2019}, and most of the northern interior has no station
  records \citep{wang2026_essd}.
\item
  \textbf{No placement study for surface melt.} Polar placement studies
  target air temperature
  \citep{hakim2020, tardif2022, kim2024, andersson2023} or sea ice
  \citep{kaminski2015, kaminski2018}. None targets surface melt or
  surface mass balance, which are spatially more heterogeneous and
  switch on and off with the seasons.
\item
  \textbf{No test of placement stability under climate change.} Existing
  studies optimise for a single climate period. Whether an optimal
  network stays optimal as the climate warms, and which sites are
  informative in both present and future climates, has not been tested.
\item
  \textbf{No comparison of traditional and data-driven criteria.}
  Ensemble-sensitivity and machine-learning placements have each been
  compared with random baselines, but, to our knowledge, not with each
  other on the same problem. TODO: confirm, and decide whether we
  include an ensemble-sensitivity baseline.
\end{enumerate}

This study addresses gaps 1 to 3 directly, through placement for
Greenland surface melt across present, near-future and long-future
windows (Q1 and Q2). TODO: gap 4, depending on the baseline decision.

\section{Data}\label{data}

\subsection{HIRHAM5 simulations}\label{hirham5-simulations}

HIRHAM5 is a regional climate model from the Danish Meteorological
Institute that combines HIRLAM dynamics with ECHAM5 physics on a 0.05
degree (about 5.5 km) rotated polar grid \citep{lucaspicher2012}. Its
output forces an offline multilayer snow model for meltwater
percolation, retention and refreezing \citep{langen2017}; see also
\citet{mottram2017}. We use HIRHAM5 runs forced by ERA-Interim, CESM2
and EC-Earth3. TODO: period, scenario and variables of each run (the
near-future runs in particular).

\subsection{PROMICE and GC-Net automatic weather
stations}\label{promice-and-gc-net-automatic-weather-stations}

The Programme for Monitoring of the Greenland Ice Sheet (PROMICE) has
run automatic weather stations (AWS) on the ice sheet since 2007
\citep{fausto2021}. The network was designed to fill a gap. Earlier
networks, chiefly the Greenland Climate Network (GC-Net), started in
1990 and extended in 1995, placed most of their stations in the
accumulation area, leaving the ablation area covered by only a handful
of sites. PROMICE sites were chosen to best complement the existing
stations while staying within range of heliports and airports
\citep{fausto2021}. All sites lie in slow-moving regions considered
relatively safe from crevasse formation \citep{fausto2026}. Siting
therefore balanced coverage against logistics and safety; it was not the
result of a formal optimisation.

The original PROMICE network had 14 AWS, deployed as pairs in seven
regions around the ice sheet (KPC, SCO, TAS, QAS, NUK, UPE and THU). In
each region, the lower (L) station sits near the ice-sheet margin in the
low-elevation ablation area, and the upper (U) station sits higher in
the ablation area, near or at the equilibrium-line altitude. Some
regions later gained a middle (M) station, and a few have a bedrock (B)
station off the ice \citep{fausto2021}. In 2021 the Geological Survey of
Denmark and Greenland (GEUS) took over GC-Net and integrated it with
PROMICE. The combined PROMICE \textbar{} GC-Net dataset covers 52 active
and historical sites, in both the ablation and the accumulation area
\citep{fausto2026}. Historical GC-Net records were reprocessed so that
they are compatible with present-day accumulation-area data
\citep{vandecrux2023}. TODO: station elevations (Table 1 of
\citet{fausto2021}; Table 3 of \citet{fausto2026}, neither extracted),
and the split between operational and decommissioned sites.

Station design follows the surface. In the ablation area, each station
is a free-standing aluminium tripod with one instrument boom, about 2.7
m above the surface. The tripod rests on the ice and sinks with it as
the surface ablates, with a battery box hung beneath the mast to lower
the centre of gravity. Sonic rangers on the tripod and on a separate
stake assembly drilled about 7 m into the ice measure surface height
change. A pressure transducer in a hose drilled typically more than 10 m
into the ice measures ablation directly \citep{fausto2021}. Since 2021,
a new tripod version (``v3'') has replaced the older one at most
ablation sites. In the accumulation area, stations are masts with two
booms, giving measurements at two levels; tripods can be carried by
helicopter, whereas the larger masts were flown in by Twin Otter
aircraft \citep{fausto2026}. Stations are serviced during visits of 2--4
h, and the remote KPC stations are visited only every 3--4 years
\citep{fausto2021, fausto2026}.

Each station measures air pressure, air temperature, relative humidity,
wind speed and direction, upward and downward shortwave and longwave
radiation, surface height, ablation (ablation-area stations), subsurface
temperature (a thermistor string to 10 m depth), station tilt and GPS
position. Derived variables include surface temperature and turbulent
heat fluxes. Measurements are taken every 10 min and hourly averages are
transmitted by Iridium; older stations sent only daily averages in
winter. The open-source pypromice software processes the data from raw
(L0) to research-ready (L3) products. The L3 product is given per site,
at hourly, daily and monthly resolution, and merges successive stations
within 4 km of one another into a single site record \citep{fausto2026}.
The dataset is a living product, updated as new data arrive, with
near-real-time access and a manually quality-controlled monthly release
(doi:10.22008/FK2/IW73UU).

Data quality varies between stations. 18 of the 26 stations analysed by
\citet{fausto2021} have at least 85\% data coverage for the surface
energy balance variables; the others suffered power failures, toppling,
snow burial or crevasse formation. Some stations were relocated, and ice
flow moves all of them over time. Known error sources include riming,
instrument burial and azimuth drift, and turbulent fluxes and rainfall
corrections are uncertain \citep{fausto2026}.

We use the network in two ways. TODO: confirm both. First, the station
locations define the context set of existing stations in the placement
experiment (Section 4.2), with values taken from the HIRHAM5 output at
those locations rather than from the observations, following
\citet{andersson2023}. The 4 km site radius is close to the HIRHAM5 grid
spacing (about 5.5 km), so each site maps to a single grid cell. Second,
the network's design (dense coverage of the ablation area at pairs of
fixed elevations, sparse coverage of the interior and the north
\citep{wang2026_essd}) is the baseline that the proposed placements are
compared against. TODO: which sites (operational in which years, given
the 2015-2016 analysis years), and whether historical GC-Net sites are
included.

\subsection{Climate windows and data
split}\label{climate-windows-and-data-split}

We analyse two windows. Within each, model training and sensor placement
years do not overlap (Table 1).

\begin{table}[t]
\caption{Model training and sensor placement years per climate window.}
\begin{tabular}{lll}
\tophline
Window & Train & Sensor placement \\
\middlehline
Present & 1990-2014 (25 yr) & 2016 \\
Near future & 2035-2059 (25 yr) & 2060 \\
Long future & 2075-2099 (25 yr) & 2100 \\
\bottomhline
\end{tabular}
\belowtable{}
\end{table}

The present window matches the years assessed in
\citet{schlager_preprint} (1990-2016). One model is trained per window
and driving dataset because regime matters more than the driving model
\citep{schlager_preprint}.

\section{Methods}\label{methods}

\subsection{Emulator}\label{emulator}

We use a convolutional conditional neural process (ConvCNP)
\citep{gordon2020} to predict daily melt. TODO: inputs, architecture,
loss, optimiser, number of seeds (planned: at least five per window),
seed selection (training and search years only).

\subsection{Sensor placement
experiment}\label{sensor-placement-experiment}

We follow the design of \citet{andersson2023}: context stations fixed at
real station locations, a regular search grid, acquisition functions
that rank candidate locations, and an analysis that adds the proposed
stations in order and tracks model performance. TODO: station set, grid
resolution, number of sensors K, acquisition functions, date sampling
for search and analysis, the cross-window regret matrix for Q2. The
sensitivity of placements to the search grid spacing is tested in
Appendix B.

\subsection{Evaluation}\label{evaluation}

We assess the trained ConvCNP model's ability to reconstruct surface
melt fields by analysing results from daily test tasks in the 2015-2016
testing periods. The context observations for surface melt tasks are
provided at the locations of the 39 PROMICE/GC-Net station active as of
June 2017, defined as those having at least one measurement in the
previous year. This task setup corresponds to the initialization used
for the sensor placement experiments and thus enables us to evaluate the
model's suitability for that application.

We report RMSE, MAE and mean bias error (MBE) of daily melt, the annual
melt total (Gt) and the melt extent (km\^{}2), for each analysis year.
TODO: units of RMSE, MAE and MBE (not labelled on the source plots), and
per-year spread relative to the window average.

\section{Results}\label{results}

\subsection{Emulator evaluation}\label{emulator-evaluation}

Table 2 gives the results for 2016, 2060 and 2100 (evaluation plots in
\texttt{figures/evaluation/}). Percentages are derived as (predicted -
true) / true.

\begin{table*}[t]
\caption{Emulator evaluation per year and driving dataset. 2060 and 2100 are only
available for CESM2 and EC-Earth3.}
\begin{tabular}{>{\raggedright\arraybackslash}p{\dimexpr(\textwidth-16\tabcolsep)*125/1000\relax}>{\raggedright\arraybackslash}p{\dimexpr(\textwidth-16\tabcolsep)*125/1000\relax}>{\raggedright\arraybackslash}p{\dimexpr(\textwidth-16\tabcolsep)*125/1000\relax}>{\raggedright\arraybackslash}p{\dimexpr(\textwidth-16\tabcolsep)*125/1000\relax}>{\raggedright\arraybackslash}p{\dimexpr(\textwidth-16\tabcolsep)*125/1000\relax}>{\raggedright\arraybackslash}p{\dimexpr(\textwidth-16\tabcolsep)*125/1000\relax}>{\raggedright\arraybackslash}p{\dimexpr(\textwidth-16\tabcolsep)*125/1000\relax}>{\raggedright\arraybackslash}p{\dimexpr(\textwidth-16\tabcolsep)*125/1000\relax}}
\tophline
Year & Driving & RMSE & MAE & MBE & Total melt true -\textgreater{} predicted (Gt) & Total bias & Melt extent bias \\
\middlehline
2016 & ERA-Interim & 1.53 & 0.33 & -0.23 & 757 -\textgreater{} 607 & -19.8\% & -3.3\% \\
2016 & CESM2 & 1.61 & 0.41 & +0.21 & 863 -\textgreater{} 1003 & +16.2\% & +28.4\% \\
2016 & EC-Earth3 & 1.38 & 0.35 & -0.02 & 899 -\textgreater{} 886 & -1.4\% & +20.0\% \\
2060 & CESM2 & 2.10 & 0.70 & +0.03 & 1850 -\textgreater{} 1872 & +1.2\% & +3.3\% \\
2060 & EC-Earth3 & 2.07 & 0.62 & +0.10 & 1663 -\textgreater{} 1726 & +3.8\% & +13.4\% \\
2100 & CESM2 & 3.67 & 1.24 & +0.81 & 3424 -\textgreater{} 3956 & +15.5\% & +8.8\% \\
2100 & EC-Earth3 & 3.57 & 1.26 & +0.87 & 2787 -\textgreater{} 3359 & +20.5\% & +14.4\% \\
\bottomhline
\end{tabular}
\belowtable{}
\end{table*}

Observations from Table 2 and the maps:

\begin{itemize}
\tightlist
\item
  Errors grow with the year for both driving datasets. RMSE rises from
  1.61 (2016) to 2.10 (2060) and 3.67 (2100) for CESM2, and from 1.38 to
  2.07 and 3.57 for EC-Earth3. TODO: check against melt magnitude, which
  also rises strongly (true totals of 863 Gt in 2016 and 3424 Gt in 2100
  for CESM2).
\item
  In both future years the emulator over-predicts melt for both driving
  datasets (positive MBE and total bias), and the over-prediction grows
  sharply by 2100 (+15.5\% for CESM2, +20.5\% for EC-Earth3).
\item
  CESM2 has the lowest bias and extent error in both future years;
  EC-Earth3 has a slightly lower RMSE in both future years and a lower
  MAE in 2060.
\item
  The residual maps show over-prediction concentrated along the
  ice-sheet margins, strongest along the northern and eastern margins
  for EC-Earth3 in 2060, and spreading across most of the margin for
  both datasets in 2100 (visual reading).
\item
  Single-year results are sensitive to how typical the year is. TODO:
  compare each analysis year with the window average melt.
\end{itemize}

\subsection{Sensor placement results}\label{sensor-placement-results}

TODO: Q1 (present window) and Q2 (near and long future; no-regret versus
future-only locations).

\section{Discussion}\label{discussion}

\textbf{Transferability.} The growth in error with warming is consistent
with the finding that emulators transfer poorly across climate regimes
\citep{schlager_preprint}. Options to improve transfer include richer
surface-state inputs, fine-tuning with a small amount of target-regime
data \citep{dheeshjith2024}, climate-invariant transformations
\citep{beucler2024}, and confidence-aware fallback to the physical model
\citep{heuer2025}. We test fine-tuning in Appendix A. TODO: summarise
its outcome here.

\textbf{Limitations.} The emulator is trained on and evaluated against a
regional climate model, not observations, so skill reflects agreement
with HIRHAM5 and inherits its biases. The sensor placement seeach and
analysis uses one year per window. The one-year search is short. TODO:
add further limitations as the experiments mature.

\conclusions

TODO.

\codedataavailability{TODO.}

\appendix

\section{Transfer learning across climate
windows}\label{appendix-a-transfer-learning-across-climate-windows}

This appendix tests whether a model trained on the present window can be
adapted to a future window with little future data, instead of training
a separate model per window (Section 3.3). Emulators of HIRHAM5 melt
transfer poorly from historical to future climate: relative error
roughly doubles and the annual total is biased
\citep{schlager_preprint}. Fine-tuning on a small amount of
target-regime data is one proposed remedy \citep{dheeshjith2024}.

\textbf{Setup.} We compare three models for each future window and
driving dataset (CESM2, EC-Earth3):

\begin{enumerate}
\def\labelenumi{\arabic{enumi}.}
\tightlist
\item
  \textbf{Zero-shot:} the present-window model applied unchanged to the
  future window.
\item
  \textbf{Fine-tuned:} the present-window model further trained on a
  subset of the future training years. TODO: number of years used
  (e.g.~1, 5, 10), which layers are updated, learning rate and number of
  epochs.
\item
  \textbf{Window-specific:} the model trained on all 25 future training
  years (Table 2 baseline).
\end{enumerate}

All three are evaluated on the same analysis years (2060 and 2100) with
the metrics of Section 4.3. TODO: number of seeds per configuration;
seed selection on training and search years only.

\textbf{Comparison.} We report RMSE, MAE, MBE, total melt bias and melt
extent bias for each model, and fine-tuned skill as a function of the
number of future years used. TODO: if the fine-tuned model reaches
window-specific skill, also compare the sensor placements it proposes
with those of the window-specific model.

\textbf{Results.} TODO. Table A1: metrics per model, window and driving
dataset. Figure A1: skill against the number of fine-tuning years.

\section{Sensitivity to search
resolution}\label{appendix-b-sensitivity-to-search-resolution}

This appendix tests whether the proposed placements depend on the
spacing of the search grid (Section 4.2). \citet{andersson2023} used one
query location per 100 km over Antarctica. The HIRHAM5 grid is much
finer (about 5.5 km), and surface melt varies over shorter distances
than air temperature, especially across the narrow ablation area along
the ice-sheet margin. A coarse grid may miss informative sites there,
while a fine grid makes the search more costly.

\textbf{Setup.} We repeat the present-window search (Q1) with the same
trained model, context stations, search dates and acquisition functions,
and change only the grid spacing. TODO: grid spacings tested
(e.g.~native about 5.5 km, then coarser multiples), the ice mask applied
to each grid, the number of candidate cells S per spacing, and the
number of sensors K.

\textbf{Comparison.} For each spacing we report (i) the analysis metrics
against the number of added sensors (Section 4.3), (ii) the overlap of
the top-K locations between spacings, matched within a distance
tolerance because the candidate sets differ, and (iii) computational
cost (number of forward passes and run time). TODO: distance tolerance
for matching locations; whether the comparison is repeated for the
future windows (Q2).

\textbf{Results.} TODO. Figure B1: proposed locations for each spacing.
Table B1: top-K overlap, analysis metrics and cost per spacing.

\textbf{Choice for the main experiments.} TODO: state the spacing used
in Section 4.2 and why.

\noappendix

\authorcontribution{TODO.}

\competinginterests{TODO.}

\begin{acknowledgements}
TODO.
\end{acknowledgements}

\bibliographystyle{copernicus}
\bibliography{references.bib}

\end{document}
